%!TEX program = lualatex
\documentclass[ngerman,headheight=70pt]{scrartcl}
\usepackage[ngerman]{babel}
\usepackage{blindtext}
\usepackage[utf8]{luainputenc}
\usepackage{scrlayer-scrpage}
\usepackage[T1]{fontenc}
\usepackage{fontspec}
\usepackage{multicol}
\usepackage{enumitem}
\usepackage{epigraph}
\usepackage[german=quotes]{csquotes}
\usepackage[bottom=2cm,footskip=8mm]{geometry}
\setmainfont{Source Sans Pro}
\MakeOuterQuote{"}

\setlength{\parindent}{0pt}
\setlength{\parskip}{6pt}

\pagestyle{plain.scrheadings}
\thispagestyle{scrheadings}
\lohead{\vspace*{-1.5cm}\includegraphics[scale=0.4]{../images/stupa_logo_rot.pdf}}
\rohead{\\Präsidium des\\Studierendenparlaments}

\lofoot[Universität Hamburg · Präsidium des Studierendenparlaments\\\hrule Von-Melle-Park 5 · D-20146 Hamburg]{Universität Hamburg · Präsidium des Studierendenparlaments\\\hrule Von-Melle-Park 5 · D-20146 Hamburg}
\cofoot[]{}
\cohead[\thepage]{}

\renewcommand{\thesubsection}{TOP \arabic{subsection}}
\renewcommand{\thesubsubsection}{\arabic{subsubsection}.}
\setcounter{subsection}{-1}

\begin{document}
    UHH · StuPa-Präsidium · Von-Melle-Park 5 · D-20146 Hamburg

    \section*{Einladung zur achten Sitzung des Studierendenparlaments}

    \textit{Liebe Parlamentarier*Innen,\\
    liebe AStA-Referent*Innen,\\
    liebe Interessierte,}

    hiermit laden wir euch zur achten Sitzung des Studierendenparlaments
    in der Legislatur 2016/2017 ein:

    \textbf{Donnerstag, den 27. Oktober 2016\\
    Um 18.00 Uhr c.t.\\
    Im Raum 0079\\
    Von-Melle-Park 5\\
    20146 Hamburg}

    Folgende \underline{Tagesordnung} schlagen wir vor:

    \underline{TEIL A}\\
    Top 0 Formalia (insgesamt 105 Min.)
    \begin{enumerate}
        \item Geschäftsbericht des Präsidiums
        \item Anfragen an das Präsidium
        \item Geschäftsbericht des AStA
        \item Anfragen an den AStA
        \item Dringlichkeitsanträge des AStA
        \item Aktuelle Stunde (falls entsprechender Antrag vorliegt)
        \item Feststellung der endgültigen Fassung des Teils B der Tagesordnung
        \item Feststellung der Beschlussfähigkeit
        \item Genehmigung der Protokolle der vorangegangenen Sitzungen
    \end{enumerate}
    \underline{TEIL B}\\
    \begin{enumerate}[label={\textbf{Top \theenumi}},leftmargin=*]
        \item Vorbereitung der Wahlen zum Studierendenparlament (60 Min.)
        \begin{enumerate}
            \item Diskussion über Konzeptvorschlag des Präsidiums (V1617-034)
            \item Diskussion über die Steigerung der Wahlbeteiligung
            \item Beschlussfassung über den Zeitplan und den Urnenplan
            (V1617-035, V1617-036)
        \end{enumerate}
        \item Debatte über den Haushalt der Freien und Hansestadt Hamburg
              2017/2018 (60 Min.)
        \item Verschiedenes (5 Min.)
    \end{enumerate}
\end{document}
